%% Plantilla mínima Jinja2+LaTeX para un generador propio.
%% Variables: las que devuelve datos_estudiante() en generar_mi_practica.py, más modo, primero, ultimo y semilla.
%% El bloque "primero" emite el preámbulo una sola vez; "ultimo" cierra \end{document}.
{% if primero -%}
\documentclass[11pt]{article}
\usepackage[spanish]{babel}
\usepackage[utf8]{inputenc}
\usepackage[T1]{fontenc}  %% comillas « » y acentos correctos
\usepackage{amsmath}
\usepackage{longtable, array, pdflscape}  %% hoja de respuestas al inicio de las soluciones
\usepackage{geometry}
\usepackage{fancyhdr}
\geometry{margin=2cm}
\pagestyle{fancy}
\fancyhf{}
\lhead{\small MI PR\'ACTICA}                       %% EDITAR: encabezado
\rhead{\small \today}
{% if modo == 'soluciones' %}\lfoot{\scriptsize Semilla: {{ semilla }}}{% endif %}  %% permite repetir la ejecución
\setlength{\headheight}{15pt}
\setlength{\parindent}{0pt}
\begin{document}
{{ clave_tex }}          %% vacío en las prácticas; en las soluciones, la hoja de respuestas
{%- endif %}

\begin{center}
    {\Large \textbf{Práctica de porcentajes}}\\[0.3em]   %% EDITAR: título
    {\normalsize Estudiante: \textbf{\underline{ {{ estudiante }} }} }
\end{center}

\noindent\textbf{Indicador:} Calcular e interpretar porcentajes a partir de frecuencias.   %% EDITAR: indicador

\vspace{0.3cm}

{% if modo == 'practicas' %}                       %% ---------- ENUNCIADO ----------

%% EDITAR: redacte la pregunta con las variables de datos_estudiante()
\textbf{1. } En una encuesta a {{ total }} {{ poblacion }}, {{ k }} {{ rasgo }}.
\begin{enumerate}
\item[a)] ¿Qué porcentaje de los encuestados {{ rasgo }}?
\item[b)] ¿Qué porcentaje no cumple esa condición?
\item[c)] Redacte una frase que comunique el resultado a la comunidad.
\end{enumerate}

{% else %}                                         %% ---------- SOLUCIÓN ----------

\textbf{1. Solución}
\begin{enumerate}
\item[a)] $\dfrac{ {{ k }} }{ {{ total }} } \cdot 100 {{ signo }} {{ porcentaje_tex }}\,\%$
\item[b)] $100 - {{ porcentaje_tex }} {{ signo }} {{ complemento_tex }}\,\%$
\item[c)] Respuesta abierta. Por ejemplo: «El {{ porcentaje }}\,\% de los {{ poblacion }} encuestados {{ rasgo }}».
\end{enumerate}

{% endif %}

{% if not ultimo %}\newpage{% endif %}             %% salto entre estudiantes, no después del último
{% if ultimo %}
\end{document}
{% endif %}
